% !TeX program = xelatex
\documentclass[11pt,letterpaper]{scrartcl}

\usepackage{fontspec}
\defaultfontfeatures{Ligatures=TeX}
\usepackage{polyglossia}
\setmainlanguage{french}
\usepackage[margin=2.5cm]{geometry}
\usepackage{amsmath,amssymb}
\usepackage{booktabs,array}
\usepackage{enumitem}
\usepackage{xcolor}
\usepackage[most]{tcolorbox}
\usepackage{listings}
\usepackage[colorlinks=true,linkcolor=accent,urlcolor=accent]{hyperref}

\definecolor{accent}{HTML}{0E9E6C}
\setlist{itemsep=2pt,topsep=4pt}
\setlength{\parindent}{0pt}
\setlength{\parskip}{0.6em}

% Environnements (mêmes couleurs que le livre)
\newtcolorbox{envremarque}{enhanced,breakable,colback=gray!8,colframe=gray!8,
  borderline west={3pt}{0pt}{gray!70},boxrule=0pt,arc=2pt,left=10pt}
\newtcolorbox{envexemple}{enhanced,breakable,colback=violet!6,colframe=violet!6,
  borderline west={3pt}{0pt}{violet!70},boxrule=0pt,arc=2pt,left=10pt}
\newtcolorbox{envdefinition}{enhanced,breakable,colback=blue!5,colframe=blue!5,
  borderline west={3pt}{0pt}{blue!70!black},boxrule=0pt,arc=2pt,left=10pt}

\lstset{language=R,basicstyle=\ttfamily\small,backgroundcolor=\color{gray!6},
  frame=single,rulecolor=\color{gray!30},breaklines=true,columns=fullflexible,
  keywordstyle=\color{blue!60!black},commentstyle=\color{gray},
  showstringspaces=false,xleftmargin=4pt,xrightmargin=4pt}

\addtokomafont{section}{\color{accent}}
\addtokomafont{subsection}{\color{accent}}

\title{Projet de session et laboratoires\\\large STT-4300 -- Régression II}
\author{Julien Miron}
\date{Hiver 2027}

\begin{document}
\maketitle
\tableofcontents
\bigskip

%=============================================================
\section{Projet de session}

\subsection{Vue d'ensemble d}

Le projet de session consiste à mener, en équipe, une analyse complète d'un jeu de données à l'aide d'un GLM : une régression logistique (réponse binaire ou proportion) ou une régression de Poisson (réponse de comptage). Le choix du modèle et du jeu de données revient à votre équipe, sous réserve de mon approbation au Labo~1.

Le projet avance par étapes, au rythme des cinq laboratoires faits en classe. Chaque laboratoire produit un livrable court qui s'intègre à votre rapport final, de sorte qu'il n'y a pas de gros bloc de travail à la fin de la session.

\begin{envremarque}
Objectif général : passer d'une question concrète à un modèle défendable, c'est-à-dire un modèle dont vous savez expliquer le choix, vérifier l'ajustement, interpréter les coefficients et reconnaître les limites.
\end{envremarque}

\subsection{Modalités}

\begin{itemize}
  \item Équipes de 3 ou 4 étudiants, formées au Labo~1.
  \item Un dépôt GitHub par équipe, avec un projet R reproductible (\texttt{renv} recommandé). Je dois avoir accès au dépôt. kkk
  \item Tout le code est en R et le rapport est rédigé en Quarto ou en \LaTeX{} : le document doit se compiler de zéro sans intervention manuelle.
\end{itemize}

\subsection{Choix du jeu de données}

Votre jeu de données doit respecter les critères suivants.
\begin{enumerate}
  \item Une variable réponse $Y$ claire : binaire (ou proportion $y_i/n_i$) pour la régression logistique, ou comptage pour la régression de Poisson.
  \item Au moins $n = 200$ observations et au moins 6 covariables potentielles, avec un mélange de variables quantitatives et qualitatives.
  \item Une question de recherche que vous pouvez énoncer en une phrase.
  \item Des données libres de droits et sans information permettant d'identifier des personnes.
  \item Un peu de désordre (valeurs manquantes, catégories rares, valeurs aberrantes) : le nettoyage fait partie de l'exercice.
\end{enumerate}

Trouvez votre jeu de données, par exemple à partir de Données Québec, de Statistique Canada, du dépôt UCI, de Kaggle ou de données ouvertes de la Ville de Québec (collisions, 311, bacs de recyclage, etc.).

\subsection{Livrable final}

Le projet se termine par une présentation en classe (10 minutes plus 5 minutes de questions) accompagnée d'un produit de communication, au choix de votre équipe :
\begin{itemize}
  \item une application Shiny qui permet d'explorer le modèle (prédictions selon les valeurs des covariables, effet d'une variable, diagnostics) ;
  \item un site web Quarto (ou un tableau de bord) racontant l'analyse ;
  \item un poster scientifique (format $36 \times 48$ po) ;
  \item un autre format du même niveau, à me proposer (balado, bande dessinée statistique, vidéo de 5 minutes, etc.).
\end{itemize}

Peu importe le format, vous devez aussi remettre un rapport technique de 8 à 12 pages qui rassemble les livrables des laboratoires et la conclusion.

Le produit de communication s'adresse à un public non statisticien. Votre présentation doit donc traduire les résultats en langage clair : parlez de rapport de cotes ou de variation en pourcentage du taux, pas seulement de coefficients $\hat\beta_j$.

\subsection{Contenu attendu du rapport}

\begin{enumerate}
  \item Question de recherche et description des données.
  \item Nettoyage et analyse exploratoire.
  \item Modèle retenu, avec sa justification.
  \item Estimation, intervalles de confiance et tests (valeur-$p$) pour les paramètres d'intérêt.
  \item Sélection de variables et comparaison de modèles.
  \item Diagnostics et validation (résidus, observations influentes, surdispersion pour Poisson, calibration et courbe ROC pour la logistique, validation croisée).
  \item Interprétation, limites et ouverture vers d'autres modèles.
  \item Une section « Ce qui pourrait fausser nos conclusions » (biais de sélection, variables omises, différence entre corrélation et causalité). Si les données concernent des personnes, ajoutez une courte réflexion éthique : qui est dans les données et qui n'y est pas ?
\end{enumerate}

%=============================================================
\section{Laboratoires}

Les cinq laboratoires ont lieu en classe, pendant environ une heure chacun. Je circule entre les équipes pour répondre à vos questions. Chaque laboratoire se termine par un livrable court (une à deux pages ou un fichier source) à remettre dans la semaine qui suit.

\subsection{Labo 1 : choix du jeu de données}

Objectif : former les équipes, choisir une question de recherche et valider la faisabilité.

Tâches en classe :
\begin{itemize}
  \item explorer au moins deux jeux de données candidats avec \texttt{str()}, \texttt{summary()} et \texttt{skimr::skim()} ;
  \item énoncer une question de recherche et identifier la variable réponse $Y$ ainsi que les covariables potentielles ;
  \item justifier le choix entre logistique et Poisson ;
  \item vérifier les critères du choix de jeu de données ci-dessus.
\end{itemize}

\begin{envexemple}
Pour un jeu de données de collisions routières par intersection, vous pourriez poser $Y_i \sim \text{Poisson}(\mu_i)$ avec
\[
  \log(\mu_i) = \log(t_i) + \beta_0 + \beta_1 x_{i1} + \cdots + \beta_p x_{ip},
\]
où $t_i$ est une mesure d'exposition (ici, le trafic moyen) qui entre dans le modèle comme décalage (\texttt{offset}).
\end{envexemple}

Livrable : une page de proposition (question, données, variable réponse, modèle envisagé, répartition des tâches). Je l'approuve ou je demande des ajustements.

\subsection{Labo 2 : nettoyage et analyse primaire}

Objectif : connaître ses données avant de modéliser.

Tâches en classe :
\begin{itemize}
  \item documenter les variables (type, unité, valeurs permises) dans un dictionnaire de données ;
  \item traiter les valeurs manquantes (quantifier, comprendre le mécanisme, choisir une stratégie et la justifier) ;
  \item recoder les variables qualitatives (niveaux rares, catégorie de référence) ;
  \item produire des statistiques descriptives et des graphiques de $Y$ en fonction de chaque covariable ;
  \item repérer les valeurs aberrantes et les problèmes de colinéarité (corrélations, facteurs d'inflation de la variance) ;
  \item mettre de côté au hasard $20\,\%$ des observations (jeu de réserve) sans y toucher avant le Labo~4 : elles serviront à évaluer honnêtement votre modèle final.
\end{itemize}

\begin{lstlisting}
library(tidyverse)
donnees <- read_csv("donnees.csv")

# proportion de Y selon une covariable qualitative
donnees |>
  group_by(covariable) |>
  summarise(n = n(), prop = mean(y), .groups = "drop")

# valeurs manquantes
donnees |> summarise(across(everything(), ~ mean(is.na(.x))))
\end{lstlisting}

Livrable : un fichier source qui part des données brutes et produit les données nettoyées et le jeu de réserve, avec une section expliquant chaque décision prise.

\subsection{Labo 3 : ajustement et début d'inférence}

Objectif : ajuster un premier GLM et interpréter les résultats.

Tâches en classe :
\begin{itemize}
  \item ajuster un modèle complet avec \texttt{glm()} ;
  \item interpréter chaque coefficient sur l'échelle du prédicteur linéaire et sur l'échelle naturelle (rapport de cotes $e^{\hat\beta_j}$ ou rapport de taux) ;
  \item construire des intervalles de confiance de Wald et par vraisemblance profilée (\texttt{confint()}), puis comparer ;
  \item effectuer des tests de Wald et des tests du rapport de vraisemblance, en précisant la valeur-$p$ et l'hypothèse testée ;
  \item tracer l'effet d'une covariable sur la probabilité (ou le taux) prédite, les autres covariables étant fixées.
\end{itemize}

\begin{lstlisting}
ajust <- glm(y ~ ., data = donnees_propres, family = binomial)   # ou poisson
summary(ajust)
exp(cbind(OR = coef(ajust), confint(ajust)))
anova(ajust, test = "Chisq")
\end{lstlisting}

\begin{envremarque}
Pour Poisson, vérifiez dès cette étape la surdispersion en comparant la déviance résiduelle à ses degrés de liberté. Si le rapport est nettement supérieur à 1, passez à un modèle quasi-Poisson ou binomial négatif.
\end{envremarque}

Livrable : une page d'interprétation écrite pour un lecteur non statisticien.

\subsection{Labo 4 : sélection de variables}

Objectif : comparer des modèles de façon rigoureuse.

Tâches en classe :
\begin{itemize}
  \item sélection pas à pas (\texttt{step()}) avec l'AIC et le BIC, et discussion des limites de cette approche, notamment l'inférence faite après sélection ;
  \item régularisation ridge et lasso avec \texttt{glmnet}, choix de $\lambda$ par validation croisée ;
  \item comparaison du pouvoir prédictif par validation croisée, puis une évaluation finale du modèle retenu sur le jeu de réserve (taux de bonne classification, AUC et déviance pour la logistique ; erreur quadratique et déviance pour Poisson) ;
  \item examen de la stabilité : les variables retenues changent-elles si l'on rééchantillonne (bootstrap) ?
\end{itemize}

\begin{lstlisting}
library(glmnet)
X <- model.matrix(y ~ ., donnees_propres)[, -1]
cv <- cv.glmnet(X, donnees_propres$y, family = "binomial", alpha = 1)
coef(cv, s = "lambda.1se")
\end{lstlisting}

Livrable : un tableau comparatif d'au moins trois modèles (complet, sélectionné, lasso), avec un choix final justifié.

\subsection{Labo 5 : ouverture sur d'autres modèles}

Objectif : voir jusqu'où le GLM est adéquat et ce qui pourrait le remplacer ou le compléter.

Tâches en classe (choisissez-en au moins une) :
\begin{itemize}
  \item ajouter une relation non linéaire par splines ou par un GAM (\texttt{mgcv::gam()}) et comparer avec le GLM au moyen de l'AIC et de la validation croisée ;
  \item essayer une méthode de plus proches voisins ou à noyau pour la prédiction ;
  \item tester un modèle pour données surdispersées (quasi-Poisson, binomial négatif) ou avec excès de zéros (modèle à inflation de zéros) ;
  \item comparer avec une méthode d'apprentissage automatique (forêt aléatoire, arbre de décision) et discuter du compromis entre interprétabilité et prédiction.
\end{itemize}

\begin{lstlisting}
library(mgcv)
gam1 <- gam(y ~ s(diametre) + s(hauteur) + espece, data = donnees_propres,
            family = binomial, method = "REML")
plot(gam1, pages = 1)
AIC(ajust, gam1)
\end{lstlisting}

Livrable : une demi-page répondant à la question « le GLM suffit-il ? », appuyée par au moins une comparaison chiffrée.

%=============================================================
\section{Évaluation}

Le projet compte pour $45\,\%$ de la note finale.

\begin{center}
\begin{tabular}{@{}lc@{}}
\toprule
Composante & Pondération \\
\midrule
Livrables des laboratoires ($5 \times 3\,\%$) & $15\,\%$ \\
Rapport technique & $15\,\%$ \\
Présentation et produit de communication & $10\,\%$ \\
Contribution individuelle (évaluation par les pairs) & $5\,\%$ \\
\midrule
Total & $45\,\%$ \\
\bottomrule
\end{tabular}
\end{center}

Critères d'évaluation du rapport et de la présentation :
\begin{itemize}
  \item justesse statistique (modèle adapté, hypothèses vérifiées, inférence correcte) ;
  \item qualité de l'interprétation (claire, nuancée, liée à la question de départ) ;
  \item reproductibilité (le document se compile sans intervention) ;
  \item clarté de la communication (graphiques lisibles, vulgarisation) ;
  \item profondeur de l'analyse (diagnostics, sélection de modèle, ouverture).
\end{itemize}

%=============================================================
\section{Calendrier suggéré}

\begin{center}
\begin{tabular}{@{}cll@{}}
\toprule
Semaine & Matière en cours & Activité \\
\midrule
3  & Régression logistique & Labo 1 : choix du jeu de données \\
5  & GLM, Poisson & Labo 2 : nettoyage et analyse primaire \\
7  & GLM, inférence & Labo 3 : ajustement et inférence \\
10 & Validation croisée, régularisation & Labo 4 : sélection de variables \\
12 & Splines, GAM & Labo 5 : ouverture \\
14 & & Présentations \\
\bottomrule
\end{tabular}
\end{center}

\end{document}
